\documentclass[11pt]{article}
\usepackage{mzqastyle}
\setrunningtitle{AI Analyst — System \& Modeling Reference}

\begin{document}

\mzqatitlepage
  {Comprehensive single-volume reference}
  {AI Analyst}
  {System \& Modeling Reference}
  {AI Analyst is a multi-agent equity-research application: a LangGraph investment
   committee of LLM analysts deliberating over a deterministic valuation engine, a
   value/sentiment screener, and a quantitative desk that adds machine-learned
   return and risk forecasting and portfolio optimization built on Microsoft
   \code{qlib}, all on top of a standardized, reconciled fundamentals warehouse.

   \vspace{8pt}
   This is the comprehensive single-volume reference. It describes the system and
   gives the essential specification of every modeling approach --- the
   gradient-boosted alpha model, the factor-structured risk models, the portfolio
   optimizers, the walk-forward backtest, the rolling factor regression, the
   macro-regime models, and the deterministic valuation and scoring models. Full
   derivations, algorithms, and hyperparameters live in the companion volumes:
   \emph{System Architecture \& Operations}, \emph{Quantitative \& ML Models}, and
   \emph{SEC \& EDINET Metrics Mapping \& Reconciliation}.}
  {System \& Modeling Reference \quad·\quad Generated 31 July 2026 \quad·\quad Independent research, not investment advice.}

\mzqatoc
\clearpage

%======================================================================
\part*{Part I\quad System}
\addcontentsline{toc}{section}{Part I\quad System}

\section{Overview and architecture}
AI Analyst turns a question --- \emph{``what is this stock worth, and what should I
expect from it?''} --- into a structured, evidence-anchored answer. It has three
qualitative entry modes (single-stock committee, group/industry relative-value
verdict, and an idea scanner) and a quantitative \emph{Quant} desk.

\begin{center}\small
\begin{tabular}{@{}>{\bfseries\color{navy}}l l@{}}
\toprule
\tabhead{Layer} & \tabhead{Technology} \\
\midrule
Backend          & FastAPI (async), Uvicorn, port 8027; blocking work on worker threads \\
Committee engine & LangGraph \code{StateGraph}; custom LLM-calling nodes \\
Quant layer      & \code{qlib} (LightGBM, scikit-learn, cvxpy) --- lean in-process embed \\
Frontend         & Next.js 14 / React 18 / Tailwind, port 3027 (REST only) \\
Database         & read-only PostgreSQL \code{xbrl\_sec} warehouse (schema \code{sec}) \\
LLM layer        & 5 providers (DeepSeek, OpenAI, Anthropic, Moonshot, Gemini), per request \\
\bottomrule
\end{tabular}
\end{center}

Three backend packages carry the logic: \code{ai\_analyst/committee} (the tribunal
graph, valuation, and evidence), \code{api} (REST routers and the \code{api/quant}
modeling layer), and \code{xbrl\_sec} (data access, the macro-cycle models, and the
LLM client). The browser is a thin client; all financial logic and every LLM call
live in the backend, whose only persistent dependency is the shared warehouse.
The service topology, the committee graph, and setup are documented in full in the
\emph{System Architecture \& Operations} volume.

\section{The data warehouse and standardized layer}
A pre-built, read-only warehouse is the single source of truth; the application
never writes to it. The fundamentals it serves are not raw filings but a
\emph{standardized} layer: raw XBRL facts from SEC EDGAR (US) and EDINET (JP) are
mapped onto one governed vocabulary of standardized line items, derived where
filers omit subtotals, checked against accounting identities, turned into a metric
library, and reconciled back to their source concepts and filings. Fundamentals are
aligned point-in-time --- a value with fiscal-period end $p$ becomes usable only at
$\operatorname{monthend}(p+90\text{d})$ --- so a signal never uses data that was not
yet public. Key tables: prices (\code{fact\_prices\_\{us,jp,etf\}}), fundamentals
(\code{fact\_metrics\_\{us,jp\}}, \code{fact\_fundamentals\_std\_\{us,jp\}}), factors
(\code{fact\_fama\_french}, \code{fact\_factor\_loadings}), macro regime
(\code{fact\_cycle\_state\_monthly}, \code{fact\_equity\_factor\_ic\_regime}),
text/sentiment (\code{fact\_mda\_sections\_\{us,jp\}}), and 13F ownership. The
mapping, derivation, metrics, and reconciliation are specified in the \emph{SEC \&
EDINET Metrics Mapping \& Reconciliation} volume.

\section{The investment committee}
A LangGraph state machine over an \code{InvestmentCommitteeState}. It splits into a
\emph{prepare} phase (a deterministic data gate $\to$ the financial-analysis engine
$\to$ evidence nodes: news/macro, 13F, a data-quality agent, and qlib quant
signals) and a \emph{debate} phase (Advocate, Challenger, Auditor and sector
specialists $\to$ Lead $\to$ memo). The prepare phase is provider-independent, so
several LLM providers can debate the same prepared evidence concurrently.
Specialists include a Quantitative Factor Analyst that interprets the quant-desk
signals. The output is a probability-weighted fair value, a triangulation range,
scenarios, an evidence bundle, and a written memo --- a valuation with its
reasoning, not a single buy/sell label. A deterministic completeness gate is a hard
stop; the accounting-identity data-quality check is advisory by default and blocks
only under an opt-in strict mode.

\section{Screener and idea generation}
The screener filters the US/JP universe on fundamentals (deterministic SQL) and
offers a natural-language screen (LLM $\to$ filters). The \emph{value $+$ sentiment}
scanner produces a $0$--$100$ interest score (\S\ref{sec:interest}); the group
ranking produces a deterministic relative-value composite (\S\ref{sec:group}). The
prompt screen fans out across every keyed provider and unions the results; the
group verdict runs a single structured narrative over the deterministic ranking.

\section{The quant desk}
A lean, in-process integration of \code{qlib}: a cross-sectional alpha model
(return prediction, \S\ref{sec:alpha}), a factor-structured risk model
(\S\ref{sec:risk}), dual portfolio optimization (native \emph{and} qlib backends,
\S\ref{sec:opt}), and a walk-forward out-of-sample backtest with Fama--French
benchmarking (\S\ref{sec:bt}). It is exposed via a \code{/api/quant/*} router and a
\emph{Quant} tab, and fed into the committee as an evidence node. The engineering
of the embed is documented in the \emph{qlib Integration} volume.

\clearpage
%======================================================================
\part*{Part II\quad Modeling reference}
\addcontentsline{toc}{section}{Part II\quad Modeling reference}

This part gives the essential specification of each model --- its object, its
estimator, and its key equations. Full derivations, algorithms, and shipped
hyperparameters are in the \emph{Quantitative \& ML Models} volume.

\section{Notation}
$\wv\in\R^N$ portfolio weights; $\muv\in\R^N$ expected (annualized) returns;
$\Sigma\in\R^{N\times N}$ return covariance; $r_f$ risk-free rate; $\ones$ the ones
vector; $z$ a cross-sectional $z$-score. Time subscripts $t$ index months unless
stated; annualization uses $252$ trading days (daily) or $12$ (monthly).

\section{Cross-sectional alpha model}
\label{sec:alpha}
\code{api/quant/qlib\_alpha.py}, on qlib's \code{LGBModel}. On a monthly,
cross-sectionally $z$-scored, point-in-time fundamentals panel, a gradient-boosted
tree ensemble $F_M(\zv)=\sum_{m=1}^{M}\eta\,h_m(\zv)$ predicts the forward 1m/3m
return. Each stage fits the residual (the negative gradient of the squared-error
loss), with LightGBM's leaf-wise histogram split gain
\begin{equation}
  \text{Gain}=\tfrac12\!\left[
  \frac{(\sum_{I_L}g_i)^2}{\sum_{I_L}h_i+\lambda}
  +\frac{(\sum_{I_R}g_i)^2}{\sum_{I_R}h_i+\lambda}
  -\frac{(\sum_{I}g_i)^2}{\sum_{I}h_i+\lambda}\right]-\gamma .
\end{equation}
Predictive power is the per-month information coefficient
$\IC_t=\Corr(\hat y_{\cdot,t},y_{\cdot,t})$ and its rank analogue $\RankIC_t$,
summarized by $\mathrm{ICIR}=\overline{\IC}/\Std(\IC)$. Performance is quoted from
the walk-forward out-of-sample path (\S\ref{sec:bt}), which is materially weaker
than the in-sample test segment.

\section{Factor-structured risk models}
\label{sec:risk}
\code{api/quant/qlib\_risk.py}, \code{api/quant/risk.py}. Returns are driven by
$K$ latent factors, giving a low-rank-plus-diagonal covariance
\begin{equation}
  \Sigma \;=\; F\,\Sigma_B\,F\tr+\diag(\mathbf d),
\end{equation}
with $F$ the factor exposures (PCA/factor analysis), $\Sigma_B$ the factor-return
covariance, and $\mathbf d$ the specific variance --- a stable shrinkage of the
noisy sample covariance, annualized $\times252$. Two alternatives are selectable:
Ledoit--Wolf shrinkage toward a constant-correlation target
$\Sigma^\star=\delta F+(1-\delta)S$, and a Fama--French \emph{fundamental}
covariance $\Sigma=B\,F_{\text{ann}}\,B\tr+D$ built from the named-factor loadings
of \S\ref{sec:ffreg}. Forward volatility per name is $\sigma_i=\sqrt{\Sigma_{ii}}$.

\section{Portfolio optimization}
\label{sec:opt}
\code{api/quant/qlib\_optimize.py} dispatches at runtime to a native optimizer or
any qlib optimizer; both consume the same $\muv,\Sigma$. The native optimizer
maximizes a multi-term mean--variance utility
\begin{equation}
  \max_{\wv}\;\wv\tr\muv-\tfrac12\lambda\,\wv\tr\Sigma\wv
  -\gamma_f\lVert S\odot a\odot(B\tr\wv-\mathbf t)\rVert_2^2
  -\gamma_{\text{to}}\lVert\wv-\wv_0\rVert_1-\gamma_c\lVert\wv\rVert_2^2
\end{equation}
subject to full investment, per-name bounds, a gross cap, sector caps, an
annualized volatility cap, and factor caps/ranges --- solved by SLSQP with an
analytic gradient, or as a cvxpy mixed-integer program when a position-count
(cardinality) constraint is set. The qlib backends provide GMV (min variance), MVO,
risk parity, inverse-vol, and enhanced indexing (benchmark-relative tracking-error,
consuming the $(F,\Sigma_B,\mathbf d)$ decomposition). For a solution $\wv^\star$ the
desk reports $\E[r]_{\text{ann}}=\wv^{\star\mathsf T}\muv$,
$\sigma_{\text{ann}}=\sqrt{\wv^{\star\mathsf T}\Sigma\wv^\star}$, and the Sharpe ratio.

\section{Walk-forward out-of-sample backtest}
\label{sec:bt}
\code{api/quant/qlib\_backtest.py}. For each test month, a booster trained on all
prior months ranks that month's names (investable universe, winsorized forward
returns), and the equal-weight top-$k$ (optionally long/short) book realizes the
next return --- out-of-sample by construction. The monthly series is scored for
annualized return, volatility, Sharpe, Sortino, max drawdown, and hit rate,
benchmarked against the Fama--French market ($\text{excess}$, tracking error,
information ratio, $\beta$), and regressed on the FF five factors plus momentum to
separate alpha from factor exposure:
\begin{equation}
  r_t-r_{f,t}=\alpha+\textstyle\sum_k\beta_k f_{k,t}+\varepsilon_t,\qquad
  t_\alpha=\hat\alpha/\mathrm{se}(\hat\alpha).
\end{equation}

\section{Rolling Fama--French factor regression}
\label{sec:ffreg}
\code{xbrl\_sec/sec/sources/factor\_model.py}. Per ticker, over a rolling $252$-day
window, excess returns are regressed on the FF factors (FF3/FF4/FF5/FF6),
$r_{i,t}-r_{f,t}=\alpha_i+\sum_k\beta_{i,k}f_{k,t}+\varepsilon_{i,t}$, estimated with
the robust Huber loss (down-weighting outlier days) and Newey--West HAC standard
errors (robust to autocorrelation/heteroskedasticity). Stored diagnostics include
adjusted $R^2$, residual volatility $\hat\sigma_\varepsilon\sqrt{252}$ (the
idiosyncratic $D$ of the FF covariance), Durbin--Watson, the condition number, and
$F$-statistics. These loadings feed the WACC (\S\ref{sec:wacc}) and the fundamental
covariance.

\section{Macro-cycle regime models}
\label{sec:regime}
\code{xbrl\_sec/sec/cycle/}. A PCA/dynamic-factor baseline compresses a monthly
multimodal feature set (macro, market, fundamental) into latent cycle factors
$\xv_t$. Two regime models run on top. A Gaussian Hidden Markov Model
($\xv_t\mid s_t{=}k\sim\mathcal N(\boldsymbol\mu_k,\Sigma_k)$, transition $A$, prior
$\pi$) fit by Baum--Welch yields the smoothed state posterior $\gamma_t(k)$; an
optional temporal multimodal VAE maximizes the evidence lower bound
$\E_{q_\phi}[\log p_\theta(\xv\mid\zv)]-\beta\,\mathrm{KL}(q_\phi\Vert p)$ and
thresholds its reconstruction-stress percentile into phases. States are mapped to
business-cycle phases (contraction / late-cycle / mid- / early-expansion),
NBER-calibrated for the US. The regime is linked to forecasting through the
regime-conditioned factor IC:
\begin{equation}
  \IC_{\text{regime}}(m)=\rho_{\text{Spearman}}\big(x_{m,i},\,r^{\text{fwd}}_i \mid \text{regime}\big),
\end{equation}
stored to \code{fact\_equity\_factor\_ic\_regime} and used to re-weight the scanner's
factor families (\S\ref{sec:interest}).

\section{Valuation models}
\label{sec:val}
\code{ai\_analyst/}. Deterministic, on a 7-year explicit horizon. A DCF projects
revenue, EBIT, NOPAT, and free cash flow, with Gordon terminal value and the
enterprise/equity bridge
\begin{equation}
  \text{TV}=\frac{\text{FCF}_7(1+g_\infty)}{\text{WACC}-g_\infty},\qquad
  \text{EV}=\sum_{i=1}^{7}\frac{\text{FCF}_i}{(1+\text{WACC})^i}+\frac{\text{TV}}{(1+\text{WACC})^7}.
\end{equation}
The cost of capital (\S\ref{sec:wacc}) builds cost of equity from the FF betas and
long-run premia, $R_e^{\text{CAPM}}=r_f+\beta_{\text{Mkt}}\text{ERP}$, blends a
synthetic-rating cost of debt, and weights them by market-value capital structure.
A reverse DCF solves for the growth or margin the market is pricing in;
sum-of-the-parts values each segment on its own growth/margin and a segment WACC;
scenario weighting and triangulation combine the DCF, SOTP, and peer-multiple ranges
into a probability-weighted fair value and a football field.

\subsection{Cost of capital}
\label{sec:wacc}
$\text{WACC}=\tfrac{E}{V}R_e+\tfrac{D}{V}R_d(1-\tau)$, with $E$ market cap, $D$ total
financial debt, $R_e$ the FF/CAPM cost of equity, and $R_d=r_f+\text{spread}$ from a
Damodaran-style synthetic rating on interest coverage.

\section{Scoring and ranking}
\subsection{Value $+$ sentiment interest score}
\label{sec:interest}
\code{screener\_agent.py}. A $0$--$100$ score blending a normalized qlib alpha
$v_\alpha$ (dominant term, $w_\alpha\approx0.40$) with a fundamental composite whose
value/growth budget is re-split by the current regime's factor IC:
\begin{equation}
  \text{score}=w_\alpha v_\alpha+(1-w_\alpha)\textstyle\sum_k w_k v_k,\qquad
  \text{interest}=100\cdot\text{score}.
\end{equation}
With no model or IC it reduces exactly to the fixed-weight composite.

\subsection{Group relative-value composite}
\label{sec:group}
\code{committee/group.py}. A cross-sectional $z$-score $z_{k,i}$ per metric (P/E,
EV/EBITDA, P/B, FCF yield, dividend, growth, margins) and a signed weight $w_k$
(negative for cheaper-is-better multiples) give
$\text{composite}_i=\sum_k w_k z_{k,i}$, ordered descending into
attractive/fair/expensive terciles with a per-name audit trail that sums exactly to
the composite. Loss-making multiples are excluded from their population. MD\&A
management tone (an LLM read in $[-1,1]$) feeds the interest score's tone term.

%======================================================================
\clearpage
\part*{Part III\quad Reference}
\addcontentsline{toc}{section}{Part III\quad Reference}

\section{How the models connect}
The \emph{valuation} path (DCF/WACC/SOTP/reverse/triangulation) answers what a
business is worth; the \emph{quant} path (alpha $\to\muv$, structured covariance
$\to\Sigma$, optimizer $\to\wv$, backtest) answers what to expect and how to size
it; the \emph{macro-cycle} path (HMM/VAE $\to$ regime-conditioned IC) conditions the
scoring on the regime; and the \emph{committee} fuses the deterministic evidence and
the quant signals into a reasoned, auditable verdict. The full model map is in the
\emph{Quantitative \& ML Models} volume.

\section{REST API (active surface)}
Base \code{http://127.0.0.1:8027}, under \code{/api}: \code{meta}, \code{screener}
($+$\code{/ai}), \code{screener/agent/value-sentiment}, \code{ai/committee}
($+$\code{/prepare}, \code{/debate}, \code{/iterate}), \code{ai/committee/group},
\code{prices}, \code{kpis}, \code{company}, \code{fx}, \code{llm}, and the quant desk
\code{quant/\{backends,alpha,risk,optimize,backtest\}}.

\section{Running}
Backend (venv Python, from \code{backend/} so \code{qlib} resolves):
\code{uvicorn api.main:app --port 8027}; frontend \code{npm run dev} (port 3027); or
\code{start\_committee.bat}. \code{qlib} is built from the bundled clone (no PyPI
wheel for Python 3.13) and installed non-editable so \code{import qlib} resolves from
any directory. The quant desk pre-warms the alpha cross-section and the walk-forward
backtest at startup, so the first user request is served from cache.

\vfill
\begin{center}\small\color{muted}
\emph{Independent, systematic research for educational purposes. Estimates are model
outputs, not investment advice.}
\end{center}

\end{document}
